
Three research questions structure the experiments:

\textbf{RQ1:} Does the OpenML corpus exhibit sufficient proxy FAIR score variance (CV > 0.15, bimodal distribution, adequate group sizes) for matched-pairs analysis?

\textbf{RQ2:} Does higher proxy Findable FAIR score predict significantly shorter TTFR after propensity matching?

\textbf{RQ3:} Does Findable sub-criteria disaggregation provide substantially stronger signal than aggregate FAIR compliance scoring?

\subsubsection{4.1 Cohorts}

\begin{table}[htbp]
\centering
\begin{tabular}{llll}
\toprule
Cohort & N & Matched Pairs & Purpose \\
\midrule
Full corpus (proxy FAIR scoring) & 5,000 & — & FAIR variance analysis (RQ1) \\
Smoke-test (synthetic timestamps) & 200 & 35 & Mechanism validation (RQ2, RQ3) \\
\bottomrule
\end{tabular}
\end{table}

The mechanism cohort (n = 200) is a synthetic proof-of-concept cohort constructed to validate the analysis pipeline prior to production-scale execution. Real upload-date metadata was not available from the OpenML bulk API; synthetic timestamps were generated to enable TTFR computation. All survival analysis results derive from this synthetic cohort.

\subsubsection{4.2 Baselines}

\textbf{Unadjusted Analysis (No Matching):} KM analysis applied to the full corpus without propensity matching. This baseline demonstrates the confounding problem that motivates the matched design.

\textbf{Aggregate FAIR Threshold (Ablation A):} KM and Cox analysis using binary aggregate quality score ≥ 0.5 as IV. This tests whether Findable sub-criteria disaggregation adds information beyond aggregate scoring.

\subsubsection{4.3 Evaluation Metrics}

\begin{table}[htbp]
\centering
\begin{tabular}{lll}
\toprule
Metric & Threshold & Purpose \\
\midrule
Log-rank p (matched KM) & < 0.05 & Primary gate for RQ2 \\
Direction (high-FAIR < low-FAIR) & Required & Direction check \\
Cox HR & > 1.2 & Secondary gate for RQ2 \\
FAIR score CV & > 0.15 & RQ1 existence check \\
Bimodality dip test p & < 0.05 & RQ1 distribution shape \\
SMD max post-matching & < 0.1 & Covariate balance \\
\bottomrule
\end{tabular}
\end{table}

\subsubsection{4.4 Hyperparameters}

\begin{table}[htbp]
\centering
\begin{tabular}{lll}
\toprule
Parameter & Smoke-test & Production Target \\
\midrule
Caliper factor & 0.8 SD & 0.2 SD \\
Minimum matched pairs & 30 & 500 \\
Observation window & 730 days & 730 days \\
Random seed & 42 & 42 \\
\bottomrule
\end{tabular}
\end{table}

